\documentclass[12pt,a4paper]{article}
\usepackage[utf8]{inputenc}
\usepackage[T1]{fontenc}
\usepackage[french]{babel}
\usepackage{hyperref}
\hypersetup{
colorlinks=true,
linkcolor=blue,
filecolor=magenta,
urlcolor=cyan
}
\usepackage{graphicx}
\usepackage{fancyhdr}
\usepackage{geometry}
\geometry{a4paper, margin=2.5cm}
\pagestyle{fancy}
\title{Encoding of a letter from Jeanne Lanvin from 1931}
\author{Jean Labusquière}
\date{18-01-2025}

\begin{document}

% Page de titre
\begin{titlepage}
\centering
\vspace*{1cm}
{\huge\bfseries Encoding of a letter from Jeanne Lanvin from 1931\\}
\vspace{1.5cm}
{\Large Jean Labusquière\\}
\vspace{0.5cm}
{\large Document original daté du 1931-09-22\\}
\vspace{1cm}
% Informations supplémentaires
\begin{minipage}{0.9\textwidth}
\begin{center}
\large\textbf{Description du document:}\\
\vspace{0.3cm}
\normalsize Typed letter from Jeanne Lanvin's closest collaborator, Jean Labusquière, addressed to Madame Gaussen and Minister Édouard Gaussen, then stationed in Stockholm, in 1931.The letter, written on the letterhead of the Jeanne Lanvin haute couture house, provides interesting insights into the challenges faced as early as 1931 regarding the employment of individuals of foreign origin.
\end{center}
\end{minipage}
\vfill
% Compiler informations
{\large Compilé par: Thibault Sabin\\}
{\large Date de compilation: 18-01-2025\\}
\vfill
\end{titlepage}

\tableofcontents

\newpage

\section{En-tête}
Figure: Lanvin's famous mother-and-child logo designed in 1924, symbolizing her bond with her daughter Marguerite.


\begin{center}
\textbf{\large Jeanne Lanvin}
\\22, FAUBOURG ST HONORE

               
                  TEL ELYSEES
               
               17-88-17-89
               41-97-46-55
            
TELEGR. JEANANVIN-PARIS 123
\end{center}


- DEAUVILLE, AVENUE DE L'HIPPODROME (Tél: 2-86)\\
- BIARRITZ, AVENUE EDOUARD VII (Tél: 179)\\
- PARIS-PLAGE, AVENUE SAINT-JEAN (Tél: 317)\\
- CANNES, B° DE LA CROISETTE (Tél: 15-34)\\
- BARCELONE, 103, RAMBLA DE CATALUNA (Télép: 75.437 (Automatica))\\
- MADRID\\

\section{Corps du texte}

\begin{flushright}PARIS, LE 22 Septembre 1931

Madame Marthe GAUSSEN

LEGATION DE FRANCE

à STOCKHOLM

\end{flushright}

Chère Madame,

Madame Jeanne LANVIN, qui est actuellement en voyage, vient de me transmettre votre lettre du 7 Septembre dernier qui n'a pu la joindre qu'après avoir été successivement transmise à plusieurs adresses au cours de ses déplacements.En s'excusant de ne pas vous répondre elle-même, Madame LANVIN me charge de vous dire qu'elle serait extrémement heureuse de pouvoir vous être agréable en cette circonstance, d'autant plus qu'elle sait combien vous avez été aimable avec nous pendant notre dernier séjour à Stockholm, mais, je ne vous surprendrai pas en vous disant que la crise qui sévit actuellement sur toutes les affaires en général, et sur celles de la couture en particulier, nous oblige à beaucoup de modération !Vous n'ignorez sans doute pas que beaucoup de nos collègues on dû restreindre dans de très sérieuses proportions le nombre de leurs employés ou de leurs ouvrières et que nous avons actuellement beaucoup de chômage dans notre métier. Les quelques maisons qui ont heureusement eu, comme la nôtre, le privilège de pouvoir conserver tous\footnote{TOUS LES MODELES LANVIN SONT DEPOSES QUICONQUE LES COPIE S'EXPOSE AUX RIGUEURS DE LA LOI DU 14 JUILLET 1909 POUR CONTREFACON. R.C. SEINE 231.151 B. SOCIETE AU CAPITAL DE 15 000 000 FRC}leurs employés, n'en sont pas moins obligées de tenir compte de cette situation et, indépendamment de notre devoir, nous éprouvons les plus grandes difficultés avec le Ministère du Travail pour maintenir chez nous des employés d'origine étrangère, auxquels on refuse les autorisations nécessaires pour être admis dans nos maisons, ou conservés même!En de telles circonstances, il ne nous paraît donc vraiment pas possible de prendre un véritable engagement vis-à-vis de votre protégée ni de prendre la responsabilité de la faire venirà Paris sans savoir, d'une part, si nous pourrions utiliser ses services et, d'autre part, si nous serions en droit de le faire.Toutefois, si elle prenait elle-même la décision de venir en France, elle n'aurait qu'à m'écrire un mot de votre part. Je m'empresserais de la recevoir et nous ferions l'impossible pour vous être agréable. Il faut espèrer d'ailleurs que les affaires s'amélioreront et que nous aurons plus de facilité dans l'avenir !En tout cas, je reste personnellement à votre disposition pour renseigner cette jeune fille et lui réserver, si elle vient à Paris,le meilleur accueil que vous puissiez souhaiter.

\section{Formule de politesse}Voulez-vous avoir l'obligeance de présenter mes respects au Ministre en le remerciant encore de toute la cordialité et de toute la bienveillance qu'il nous a témoignées lors de notre voyage et d'agréer, Chère Madame, mes hommages très respectueux et très dévoués.

\begin{flushright}
\textit{Jean Labusquière}
\end{flushright}

\end{document}